% GENERATED FILE. Edit canonical lessons, then run npm run generate:books.
% canonical-source-hash: fnv1a64:afbdce9f9fb1b39d
% canonical-lessons: GE-C175-ausfuellen, GE-C175-ausruhen, GE-C175-backen, GE-C175-baden, GE-C175-bauen

\chapter{To fill in, To rest, To bake, To bathe, To build}
\label{ch:ge-to-fill-in-to-rest-to-bake-to-bathe-to-build}
\hlchaptermodality{175}

\begin{chapteropening}
\textbf{By the end of this chapter:} \emph{I can say to fill in, to rest, to bake, to bathe and to build.}

Complete the last lesson of chapter 175: I can say to fill in, to rest, to bake, to bathe and to build.
\end{chapteropening}

\section[ausfüllen]{ausfüllen — to fill in (a form)}

\label{lesson:GE-C175-ausfuellen}



\begin{quote}
Before the new one: say the German for to take off, then the German for to tidy up.
\end{quote}

\subsection*{You'll want to know: ausfüllen}
\textbf{ausfüllen} — \textquotedblleft{}to fill in (a form)\textquotedblright{}.

\begin{cousinweb}
\textbf{ausfüllen} is \emph{aus-}, \textquotedblleft{}out\textquotedblright{}, + \emph{füllen}, \textquotedblleft{}to fill\textquotedblright{}, the same verb as English \emph{fill} (EtymDB 2.1): to fill out a form.
\end{cousinweb}

\begin{grammarlens}[title={What you've built}]
One more word for this chapter.
\end{grammarlens}

\subsection*{Your turn}
\begin{itemize}
\raggedright
  \item \emph{Say it:} \emph{ausfüllen}
  \item \emph{Say it:} \emph{ausfüllen}, once more
  \item \emph{Say it:} say \emph{ausfüllen}
  \item \emph{Recall:} say the German for to take off, then the German for to tidy up, then say \emph{ausfüllen} again
\end{itemize}

\begin{tcolorbox}[breakable,colback=teal!4,colframe=teal!35!black,title={Before you move on}]
What does ausfüllen mean? (\textquotedblleft{}To fill in (a form)\textquotedblright{}.) Say it once more.
\end{tcolorbox}
\section[ausruhen]{ausruhen — to rest (sich ausruhen)}

\label{lesson:GE-C175-ausruhen}



\begin{quote}
Before the new one: say the German for to tidy up, then the German for to fill in.
\end{quote}

\subsection*{You'll want to know: ausruhen}
\textbf{ausruhen} — \textquotedblleft{}to rest (sich ausruhen)\textquotedblright{}.

\begin{cousinweb}
\textbf{ausruhen} is \emph{aus-}, \textquotedblleft{}out, fully\textquotedblright{}, + the verb \textquotedblleft{}to rest\textquotedblright{}, built on the noun for \textquotedblleft{}rest\textquotedblright{}, Old High German \emph{ruowa} $\leftarrow$ Proto-Germanic \emph{*rōwō} (EtymDB 2.1).
\end{cousinweb}

\begin{grammarlens}[title={What you've built}]
One more word for this chapter.
\end{grammarlens}

\subsection*{Your turn}
\begin{itemize}
\raggedright
  \item \emph{Say it:} \emph{ausruhen}
  \item \emph{Say it:} \emph{ausruhen}, once more
  \item \emph{Say it:} say \emph{ausruhen}
  \item \emph{Recall:} say the German for to tidy up, then the German for to fill in, then say \emph{ausruhen} again
\end{itemize}

\begin{tcolorbox}[breakable,colback=teal!4,colframe=teal!35!black,title={Before you move on}]
What does ausruhen mean? (\textquotedblleft{}To rest (sich ausruhen)\textquotedblright{}.) Say it once more.
\end{tcolorbox}
\section[backen]{backen — to bake}

\label{lesson:GE-C175-backen}



\begin{quote}
Before the new one: say the German for to fill in, then the German for to rest.
\end{quote}

\subsection*{You'll want to know: backen}
\textbf{backen} — \textquotedblleft{}to bake\textquotedblright{}.

\begin{cousinweb}
\textbf{backen} $\leftarrow$ Middle High German \emph{backen} $\leftarrow$ Old High German \emph{backan} — inherited, the same verb as English \emph{bake}. (EtymDB 2.1)
\end{cousinweb}

\begin{grammarlens}[title={What you've built}]
One more word for this chapter.
\end{grammarlens}

\subsection*{Your turn}
\begin{itemize}
\raggedright
  \item \emph{Say it:} \emph{backen}
  \item \emph{Say it:} \emph{backen}, once more
  \item \emph{Say it:} say \emph{backen}
  \item \emph{Recall:} say the German for to fill in, then the German for to rest, then say \emph{backen} again
\end{itemize}

\begin{tcolorbox}[breakable,colback=teal!4,colframe=teal!35!black,title={Before you move on}]
What does backen mean? (\textquotedblleft{}To bake\textquotedblright{}.) Say it once more.
\end{tcolorbox}
\section[baden]{baden — to bathe}

\label{lesson:GE-C175-baden}



\begin{quote}
Before the new one: say the German for to rest, then the German for to bake.
\end{quote}

\subsection*{You'll want to know: baden}
\textbf{baden} — \textquotedblleft{}to bathe\textquotedblright{}.

\begin{cousinweb}
\textbf{baden} $\leftarrow$ Old High German \emph{badōn} $\leftarrow$ Proto-Germanic \emph{*baþōną}, \textquotedblleft{}to bathe\textquotedblright{} — inherited, the verb beside \emph{Bad}, the same as English \emph{bathe}. (EtymDB 2.1)
\end{cousinweb}

\begin{grammarlens}[title={What you've built}]
One more word for this chapter.
\end{grammarlens}

\subsection*{Your turn}
\begin{itemize}
\raggedright
  \item \emph{Say it:} \emph{baden}
  \item \emph{Say it:} \emph{baden}, once more
  \item \emph{Say it:} say \emph{baden}
  \item \emph{Recall:} say the German for to rest, then the German for to bake, then say \emph{baden} again
\end{itemize}

\begin{tcolorbox}[breakable,colback=teal!4,colframe=teal!35!black,title={Before you move on}]
What does baden mean? (\textquotedblleft{}To bathe\textquotedblright{}.) Say it once more.
\end{tcolorbox}
\section[bauen]{bauen — to build}

\label{lesson:GE-C175-bauen}



\begin{quote}
Before the new one: say the German for to bake, then the German for to bathe.
\end{quote}

\subsection*{You'll want to know: bauen}
\textbf{bauen} — \textquotedblleft{}to build\textquotedblright{}.

\begin{cousinweb}
\textbf{bauen} $\leftarrow$ Middle High German \emph{būwen} $\leftarrow$ Old High German \emph{būan}, \textquotedblleft{}to dwell, to farm\textquotedblright{} (EtymDB 2.1) — inherited, the verb beside \emph{Bauer}. Old English \emph{būan} meant \textquotedblleft{}to dwell\textquotedblright{}.
\end{cousinweb}

\begin{grammarlens}[title={What you've built}]
One more word for this chapter.
\end{grammarlens}

\subsection*{Your turn}
\begin{itemize}
\raggedright
  \item \emph{Say it:} \emph{bauen}
  \item \emph{Say it:} \emph{bauen}, once more
  \item \emph{Say it:} say \emph{bauen}
  \item \emph{Recall:} say the German for to bake, then the German for to bathe, then say \emph{bauen} again
\end{itemize}

\begin{tcolorbox}[breakable,colback=teal!4,colframe=teal!35!black,title={Before you move on}]
What does bauen mean? (\textquotedblleft{}To build\textquotedblright{}.) Say it once more.
\end{tcolorbox}
